% year: תשפ"ג
% semester: ב
% moed: א
% note: ד"ר פיטר סמובול
% source: exams/מבחנים/תשפג/מועד א תשפג סמסטר ב.pdf
% number: 6א
% points: 20
% categories: taylor, derivatives

\begin{question}
רשמו את פולינום מקלורן מסדר 2 לפונקציה $y=f(x)$ המוגדרת בצורה סתומה: $x^2+x+e^x+2y=0$.
\end{question}

\begin{solution}
נמצא $y(0)$, $y'(0)$, $y''(0)$ בגזירה סתומה.
\begin{itemize}
  \item $x=0$: $0+0+1+2y(0)=0\Rightarrow y(0)=-\frac12$.
  \item גזירה: $2x+1+e^x+2y'=0\Rightarrow y'=-\frac{2x+1+e^x}{2}$, ולכן $y'(0)=-\frac{2}{2}=-1$.
  \item גזירה נוספת: $2+e^x+2y''=0\Rightarrow y''=-\frac{2+e^x}{2}$, ולכן $y''(0)=-\frac32$.
\end{itemize}
פולינום מקלורן מסדר 2:
\[ P_2(x)=y(0)+y'(0)x+\frac{y''(0)}{2}x^2=\boxed{-\frac12-x-\frac34x^2}. \]
(בדיקה: כאן אפשר גם לבודד $y=-\frac12(x^2+x+e^x)$ ולהציב $e^x=1+x+\frac{x^2}{2}+\dots$, ומתקבל אותו פולינום.)
\end{solution}
